\begin{definition}[Refinement in ordinary physical coordinates]
    \label{def:mpdo_physical_t_map}
    \lean{MPOTensor.PhysicalSectorFactorization.NeighboringTraceFactorization.physicalTMap}
    \leanok
    \uses{def:mpdo_physical_sector_t_map,
      def:mpdo_physical_sector_closure_coordinates}
    Transporting $\mathcal T$ through the canonical two-site and three-site
    coordinate identifications gives a channel
    $\widehat{\mathcal T}:\MN{q^2}\longrightarrow\MN{q^3}$, where all
    multiple physical indices are associated to the right.
\end{definition}

\begin{theorem}[The physical refinement map is a channel]
    \label{thm:mpdo_physical_t_map_cptp}
    \lean{MPOTensor.PhysicalSectorFactorization.NeighboringTraceFactorization.physicalTMap_isKrausCPTP}
    \leanok
    \uses{def:mpdo_physical_t_map}
    The map $\widehat{\mathcal T}$ is trace-preserving and completely
    positive.
\end{theorem}

\begin{proof}\leanok
    \uses{thm:mpdo_physical_sector_t_map_cptp}
    The coordinate identifications are unitary permutation channels.
    Composing them with a trace-preserving completely positive map preserves
    both properties.
\end{proof}

% **Local fix (zero-weight quotient):** this identity uses the completed
% zero-weight preparation. Documented in
% docs/paper-gaps/cpgsv17_mpdo_zero_weight_preparation_completion.tex.
\begin{theorem}[Refinement identity in ordinary physical coordinates]
    \label{thm:mpdo_physical_t_closure_identity}
    \lean{MPOTensor.PhysicalSectorFactorization.NeighboringTraceFactorization.physicalTMap_twoSiteClosure_eq_threeSiteClosure}
    \leanok
    \uses{def:mpdo_physical_t_map}
    For every virtual matrix $X$,
    $\widehat{\mathcal T}(\widehat{\mathcal K}_2(X))
    =\widehat{\mathcal K}_3(X)$.
\end{theorem}

\begin{proof}\leanok
    \uses{thm:mpdo_physical_sector_t_closure_identity}
    Transport the sector-coordinate identity through the inverse three-site
    coordinate identification.
\end{proof}

\begin{definition}[Physical coarse-graining]
    \label{def:mpdo_physical_s_map}
    \lean{MPOTensor.PhysicalSectorFactorization.NeighboringTraceFactorization.physicalSMap}
    \leanok
    \uses{def:mpdo_physical_sector_s_map,
      def:mpdo_physical_sector_closure_coordinates}
    Transporting $\mathcal S$ through the canonical three-site and two-site
    coordinate identifications gives a channel
    $\widehat{\mathcal S}:\MN{q^3}\longrightarrow\MN{q^2}$.
\end{definition}

\begin{theorem}[The physical coarse-graining map is a channel]
    \label{thm:mpdo_physical_s_map_cptp}
    \lean{MPOTensor.PhysicalSectorFactorization.NeighboringTraceFactorization.physicalSMap_isKrausCPTP}
    \leanok
    \uses{def:mpdo_physical_s_map}
    The map $\widehat{\mathcal S}$ is trace-preserving and completely
    positive.
\end{theorem}

\begin{proof}\leanok
    \uses{thm:mpdo_physical_sector_s_map_cptp}
    This follows by composition with the two coordinate-permutation
    channels.
\end{proof}

\begin{theorem}[Coarse-graining identity in ordinary physical coordinates]
    \label{thm:mpdo_physical_s_closure_identity}
    \lean{MPOTensor.PhysicalSectorFactorization.NeighboringTraceFactorization.physicalSMap_threeSiteClosure_eq_twoSiteClosure}
    \leanok
    \uses{def:mpdo_physical_s_map}
    For every virtual matrix $X$,
    $\widehat{\mathcal S}(\widehat{\mathcal K}_3(X))
    =\widehat{\mathcal K}_2(X)$.
\end{theorem}

\begin{proof}\leanok
    \uses{thm:mpdo_physical_sector_s_closure_identity}
    Transport the sector-coordinate identity through the inverse two-site
    coordinate identification.
\end{proof}

\begin{definition}[Localized physical refinement]
    \label{def:mpdo_localized_physical_t_map}
    \lean{MPOTensor.PhysicalSectorFactorization.NeighboringTraceFactorization.localizedPhysicalTMap}
    \leanok
    \uses{def:mpdo_physical_t_map}
    Reassociating three right-associated sites as $(12)3$, applying
    $\widehat{\mathcal T}\otimes\id$ to the first two, and
    reassociating the result back to $1(2(34))$ defines the localized
    refinement channel
    $\widehat{\mathcal T}_{(12)3}:\MN{q^3}\longrightarrow\MN{q^4}$
    on right-associated physical coordinates.
\end{definition}

\begin{theorem}[The localized physical refinement is a channel]
    \label{thm:mpdo_localized_physical_t_map_cptp}
    \lean{MPOTensor.PhysicalSectorFactorization.NeighboringTraceFactorization.localizedPhysicalTMap_isKrausCPTP}
    \leanok
    \uses{def:mpdo_localized_physical_t_map}
    The localized refinement map is trace-preserving and completely
    positive.
\end{theorem}

\begin{proof}\leanok
    \uses{thm:mpdo_physical_t_map_cptp}
    Tensoring a channel with the identity preserves complete positivity and
    trace, as do the two associativity permutations.
\end{proof}

% **Local fix (zero-weight quotient):** this identity uses the completed
% zero-weight preparation. Documented in
% docs/paper-gaps/cpgsv17_mpdo_zero_weight_preparation_completion.tex.
\begin{theorem}[Localized refinement of physical closures]
    \label{thm:mpdo_localized_physical_t_closure_identity}
    \lean{MPOTensor.PhysicalSectorFactorization.NeighboringTraceFactorization.localizedPhysicalTMap_threeSiteClosure_eq_fourSiteClosure}
    \leanok
    \uses{def:mpdo_localized_physical_t_map}
    For every virtual matrix $X$,
    $(\widehat{\mathcal T}\otimes\id)
    (\widehat{\mathcal K}_3(X))
    =\widehat{\mathcal K}_4(X)$, with the map understood in
    right-associated coordinates.
\end{theorem}

\begin{proof}\leanok
    \uses{thm:mpdo_physical_t_closure_identity}
    Fixing the last ket and bra indices turns the three-site closure into a
    two-site closure with the last tensor matrix absorbed into $X$. Apply the
    two-to-three-site refinement identity to this virtual matrix.
\end{proof}

\begin{definition}[Localized physical coarse-graining]
    \label{def:mpdo_localized_physical_s_map}
    \lean{MPOTensor.PhysicalSectorFactorization.NeighboringTraceFactorization.localizedPhysicalSMap}
    \leanok
    \uses{def:mpdo_physical_s_map}
    Reassociating four right-associated sites as $(123)4$, applying
    $\widehat{\mathcal S}\otimes\id$ to the first three, and
    reassociating the result back to $1(23)$ defines the localized
    coarse-graining channel
    $\widehat{\mathcal S}_{(123)4}:\MN{q^4}\longrightarrow\MN{q^3}$
    on right-associated physical coordinates.
\end{definition}

\begin{theorem}[The localized physical coarse-graining is a channel]
    \label{thm:mpdo_localized_physical_s_map_cptp}
    \lean{MPOTensor.PhysicalSectorFactorization.NeighboringTraceFactorization.localizedPhysicalSMap_isKrausCPTP}
    \leanok
    \uses{def:mpdo_localized_physical_s_map}
    The localized coarse-graining map is trace-preserving and completely
    positive.
\end{theorem}

\begin{proof}\leanok
    \uses{thm:mpdo_physical_s_map_cptp}
    Tensor the physical coarse-graining channel with the identity and compose
    with the two associativity permutations.
\end{proof}

\begin{theorem}[Localized coarse-graining of physical closures]
    \label{thm:mpdo_localized_physical_s_closure_identity}
    \lean{MPOTensor.PhysicalSectorFactorization.NeighboringTraceFactorization.localizedPhysicalSMap_fourSiteClosure_eq_threeSiteClosure}
    \leanok
    \uses{def:mpdo_localized_physical_s_map}
    For every virtual matrix $X$,
    $(\widehat{\mathcal S}\otimes\id)
    (\widehat{\mathcal K}_4(X))
    =\widehat{\mathcal K}_3(X)$ in right-associated coordinates.
\end{theorem}

\begin{proof}\leanok
    \uses{thm:mpdo_physical_s_closure_identity}
    Fixing the last ket and bra indices identifies the four-site closure with
    a three-site closure having the last tensor matrix absorbed into $X$.
    Apply the three-to-two-site coarse-graining identity.
\end{proof}

% **Local fix (cyclic implication labels):** the concluding proof at source
% lines 1822--1824 assigns the wrong conditions to the four propositions it
% invokes. The construction below cites Proposition C.7 and its twice-applied
% channel observation directly. Documented in
% docs/paper-gaps/cpgsv17_mpdo_theorem_4_9_implication_label.tex.
\begin{definition}[Two-step physical refinement]
    \label{def:mpdo_physical_t_squared_map}
    \lean{MPOTensor.PhysicalSectorFactorization.NeighboringTraceFactorization.physicalTSquaredMap}
    \leanok
    \uses{def:mpdo_neighboring_trace_factorization,
      def:mpdo_physical_t_map,
      def:mpdo_localized_physical_t_map}
    Let $F$ be a physical-sector factorization whose neighboring operators
    satisfy Definition~\ref{def:mpdo_neighboring_trace_factorization}, and let
    $\widehat{\mathcal K}$ be its sector-coordinate tensor. Define
    \begin{align}
      \widehat{\mathcal T}^2
      &=\widehat{\mathcal T}_{(12)3}\circ\widehat{\mathcal T},
      \notag\\
      \widehat{\mathcal T}^2
      &:\MN{q^2}\longrightarrow\MN{q^4}.
      \notag
    \end{align}
    Here the second application acts on the first two sites and leaves the
    third site unchanged. Thus the superscript denotes two successive
    overlapping applications, as in
    \cite[Appendix~C.2, lines~1821--1825]{Cirac2016MPDO_arXiv}.
\end{definition}

\begin{theorem}[The two-step physical refinement is a channel]
    \label{thm:mpdo_physical_t_squared_map_cptp}
    \lean{MPOTensor.PhysicalSectorFactorization.NeighboringTraceFactorization.physicalTSquaredMap_isKrausCPTP}
    \leanok
    \uses{def:mpdo_physical_t_squared_map}
    The map $\widehat{\mathcal T}^2$ is trace-preserving and completely
    positive.
\end{theorem}

\begin{proof}\leanok
    \uses{thm:mpdo_physical_t_map_cptp,
      thm:mpdo_localized_physical_t_map_cptp}
    It is the composition of the physical refinement channel and its
    localization on the first two sites.
\end{proof}

% **Local fix (zero-weight quotient):** this identity uses the completed
% zero-weight preparation. Documented in
% docs/paper-gaps/cpgsv17_mpdo_zero_weight_preparation_completion.tex.
\begin{theorem}[Two-step refinement of the sector-coordinate tensor]
    \label{thm:mpdo_physical_t_squared_closure_identity}
    \lean{MPOTensor.PhysicalSectorFactorization.NeighboringTraceFactorization.physicalTSquaredMap_twoSiteClosure_eq_fourSiteClosure}
    \leanok
    \uses{def:mpdo_physical_t_squared_map}
    For every virtual matrix $X$, the sector-coordinate tensor satisfies
    $\widehat{\mathcal T}^2(\widehat{\mathcal K}_2(X))
    =\widehat{\mathcal K}_4(X)$.
\end{theorem}

\begin{proof}\leanok
    \uses{thm:mpdo_physical_t_closure_identity,
      thm:mpdo_localized_physical_t_closure_identity}
    The first refinement gives $\widehat{\mathcal K}_3(X)$; applying the
    localized refinement to its first two sites gives
    $\widehat{\mathcal K}_4(X)$.
\end{proof}

\begin{definition}[Two-step physical coarse-graining]
    \label{def:mpdo_physical_s_squared_map}
    \lean{MPOTensor.PhysicalSectorFactorization.NeighboringTraceFactorization.physicalSSquaredMap}
    \leanok
    \uses{def:mpdo_physical_s_map,
      def:mpdo_localized_physical_s_map}
    Define
    \begin{align}
      \widehat{\mathcal S}^2
      &=\widehat{\mathcal S}\circ\widehat{\mathcal S}_{(123)4},
      \notag\\
      \widehat{\mathcal S}^2
      &:\MN{q^4}\longrightarrow\MN{q^2}.
      \notag
    \end{align}
    Here the first application acts on the first three sites and leaves the
    fourth site unchanged. The superscript again denotes successive
    overlapping applications.
\end{definition}

\begin{theorem}[The two-step physical coarse-graining is a channel]
    \label{thm:mpdo_physical_s_squared_map_cptp}
    \lean{MPOTensor.PhysicalSectorFactorization.NeighboringTraceFactorization.physicalSSquaredMap_isKrausCPTP}
    \leanok
    \uses{def:mpdo_physical_s_squared_map}
    The map $\widehat{\mathcal S}^2$ is trace-preserving and completely
    positive.
\end{theorem}

\begin{proof}\leanok
    \uses{thm:mpdo_physical_s_map_cptp,
      thm:mpdo_localized_physical_s_map_cptp}
    It is the composition of the localized coarse-graining channel and the
    physical coarse-graining channel.
\end{proof}

\begin{theorem}[Two-step coarse-graining of the sector-coordinate tensor]
    \label{thm:mpdo_physical_s_squared_closure_identity}
    \lean{MPOTensor.PhysicalSectorFactorization.NeighboringTraceFactorization.physicalSSquaredMap_fourSiteClosure_eq_twoSiteClosure}
    \leanok
    \uses{def:mpdo_physical_s_squared_map}
    For every virtual matrix $X$, the sector-coordinate tensor satisfies
    $\widehat{\mathcal S}^2(\widehat{\mathcal K}_4(X))
    =\widehat{\mathcal K}_2(X)$.
\end{theorem}

\begin{proof}\leanok
    \uses{thm:mpdo_localized_physical_s_closure_identity,
      thm:mpdo_physical_s_closure_identity}
    The localized coarse-graining first gives
    $\widehat{\mathcal K}_3(X)$; the second coarse-graining gives
    $\widehat{\mathcal K}_2(X)$.
\end{proof}

\begin{figure}[!b]
    \centering
    % Formula: \widehat{\mathcal T}(\widehat{\cal K}_2(X))
    % =\widehat{\cal K}_3(X), followed by
    % \widehat{\mathcal T}_{(12)3}(\widehat{\cal K}_3(X))
    % =\widehat{\cal K}_4(X); the reverse row is the corresponding pair of
    % \widehat{\mathcal S} identities.
    % Ink-to-index: every \widehat{\cal K} box has west/east virtual indices
    % and upper/lower physical indices.  Each enclosure mark groups two
    % adjacent tensor bodies into one blocked local site and introduces no
    % contraction.
    % Contractions: only adjacent east--west virtual legs are summed.  All
    % upper and lower physical legs, and both outer virtual legs, remain free.
    % Boundary signatures (virtual-west, virtual-east, physical-up,
    % physical-down): the two-, three-, and four-site panels have respectively
    % (1,1,2,2), (1,1,3,3), and (1,1,4,4).  The arrows are physical channels,
    % so their endpoint physical arities change as displayed.
    % Source verdict: Appendix C.2, lines 1821--1825, specifies the two-step
    % channels algebraically but has no dedicated figure.  The rows reproduce
    % that composition using the local tensors already defined here.
    \begin{tenkzequation}
        \tenkzkernel
        $\widehat{\mathcal T}^{\,2}:\quad$
        \begin{tenkz}[cols=2, west=open, east=open, physical=updown]
            \tn[skin=mpo]{\widehat{\cal K}} & \tn[skin=mpo]{\widehat{\cal K}}
            \tnmark[form=enclosure, label pos=s]{(1,1) .. (1,2)}
                {$\widehat{\cal K}^{[2]}$}
        \end{tenkz}
        $\quad\xrightarrow{\widehat{\mathcal T}}\quad$
        \begin{tenkz}[cols=3, west=open, east=open, physical=updown]
            \tn[skin=mpo]{\widehat{\cal K}} &
            \tn[skin=mpo]{\widehat{\cal K}} &
            \tn[skin=mpo]{\widehat{\cal K}}
        \end{tenkz}
        $\quad\xrightarrow{\widehat{\mathcal T}_{(12)3}}\quad$
        \begin{tenkz}[cols=4, west=open, east=open, physical=updown]
            \tn[skin=mpo]{\widehat{\cal K}} &
            \tn[skin=mpo]{\widehat{\cal K}} &
            \tn[skin=mpo]{\widehat{\cal K}} &
            \tn[skin=mpo]{\widehat{\cal K}}
            \tnmark[form=enclosure, label pos=s]{(1,1) .. (1,2)}
                {$\widehat{\cal K}^{[2]}$}
            \tnmark[form=enclosure, label pos=s]{(1,3) .. (1,4)}
                {$\widehat{\cal K}^{[2]}$}
        \end{tenkz}
    \end{tenkzequation}
    \par    \begin{tenkzequation}
        \tenkzkernel
        $\widehat{\mathcal S}^{\,2}:\quad$
        \begin{tenkz}[cols=4, west=open, east=open, physical=updown]
            \tn[skin=mpo]{\widehat{\cal K}} &
            \tn[skin=mpo]{\widehat{\cal K}} &
            \tn[skin=mpo]{\widehat{\cal K}} &
            \tn[skin=mpo]{\widehat{\cal K}}
            \tnmark[form=enclosure, label pos=s]{(1,1) .. (1,2)}
                {$\widehat{\cal K}^{[2]}$}
            \tnmark[form=enclosure, label pos=s]{(1,3) .. (1,4)}
                {$\widehat{\cal K}^{[2]}$}
        \end{tenkz}
        $\quad\xrightarrow{\widehat{\mathcal S}_{(123)4}}\quad$
        \begin{tenkz}[cols=3, west=open, east=open, physical=updown]
            \tn[skin=mpo]{\widehat{\cal K}} &
            \tn[skin=mpo]{\widehat{\cal K}} &
            \tn[skin=mpo]{\widehat{\cal K}}
        \end{tenkz}
        $\quad\xrightarrow{\widehat{\mathcal S}}\quad$
        \begin{tenkz}[cols=2, west=open, east=open, physical=updown]
            \tn[skin=mpo]{\widehat{\cal K}} & \tn[skin=mpo]{\widehat{\cal K}}
            \tnmark[form=enclosure, label pos=s]{(1,1) .. (1,2)}
                {$\widehat{\cal K}^{[2]}$}
        \end{tenkz}
    \end{tenkzequation}
    \caption{The superscripts in $\widehat{\mathcal T}^2$ and
    $\widehat{\mathcal S}^2$ denote successive overlapping applications.
    The enclosures identify the two-site and four-site endpoints with one
    and two sites of the blocked tensor. This is the channel construction in
    \cite[Appendix~C.2, lines~1821--1825]{Cirac2016MPDO_arXiv}.}
    \label{fig:mpdo_blocked_rfp_channels}
\end{figure}

% **Local fix (zero-weight quotient):** this theorem uses the completed
% zero-weight preparation through the two-step refinement identity. Documented
% in docs/paper-gaps/cpgsv17_mpdo_zero_weight_preparation_completion.tex.
% **Scope restriction (physical coordinates):** this theorem concerns the
% sector-coordinate tensor, not yet the original tensor. Documented in
% docs/paper-gaps/cpgsv17_mpdo_blocked_rfp_physical_transport.tex.
% **Scope restriction (single supplied factorization):** this theorem assumes
% one neighboring trace factorization of the whole tensor to which it is
% applied. Theorem 4.9(iv) supplies such data only for each BNT representative,
% and does not control unequal raw repeated-copy weights. Thus this conditional
% helper does not imply the printed (iv)=>(v), which is refuted. Documented in
% docs/paper-gaps/cpgsv17_pf_rank_one.tex.
\begin{theorem}[Blocked sector-coordinate tensor is a renormalization fixed point]
    \label{thm:mpdo_blocked_sector_coordinate_tensor_is_rfp_via_ts}
    \lean{MPOTensor.PhysicalSectorFactorization.NeighboringTraceFactorization.blockTwo_sectorCoordinateTensor_isRFPViaTS}
    \leanok
    \uses{def:is_rfp_via_ts, def:mpdo_two_site_blocking,
      def:mpdo_neighboring_trace_factorization}
    Let $F$ be a physical-sector factorization whose neighboring operators
    satisfy Definition~\ref{def:mpdo_neighboring_trace_factorization}. Let
    $\widehat{\mathcal K}^{[2]}$ be the tensor obtained by blocking two
    adjacent sites of its sector-coordinate tensor $\widehat{\mathcal K}$.
    Then $\widehat{\mathcal K}^{[2]}$ is a renormalization fixed point in the
    sense of Definition~\ref{def:is_rfp_via_ts}. More precisely, after the
    canonical identifications of one blocked site with two original sites and
    two blocked sites with four original sites, the required channels are
    $\widehat{\mathcal S}^2$ and $\widehat{\mathcal T}^2$.
\end{theorem}

\begin{proof}\leanok
    \uses{thm:mpdo_physical_t_squared_map_cptp,
      thm:mpdo_physical_t_squared_closure_identity,
      thm:mpdo_physical_s_squared_map_cptp,
      thm:mpdo_physical_s_squared_closure_identity,
      thm:mpdo_one_site_closure_two_site_blocking,
      thm:mpdo_two_site_closure_two_site_blocking}
    Decode one blocked index as a pair of physical indices and two blocked
    indices as four physical indices. The one-site and two-site closures of
    $\widehat{\mathcal K}^{[2]}$ thereby become respectively
    $\widehat{\mathcal K}_2(X)$ and $\widehat{\mathcal K}_4(X)$.
    Transporting the two channels through these permutation identifications
    preserves complete positivity and trace, and the two closure identities
    give the required equations in Definition~\ref{def:is_rfp_via_ts}.
\end{proof}

% **Local fix (zero-weight quotient):** this theorem uses the completed
% zero-weight preparation through the two-step refinement identity. Documented
% in docs/paper-gaps/cpgsv17_mpdo_zero_weight_preparation_completion.tex.
% **Local fix (cyclic implication labels):** the concluding proof at source
% lines 1822--1824 assigns the wrong conditions to the four propositions it
% invokes. The construction below cites Proposition C.7 and its twice-applied
% channel observation directly. Documented in
% docs/paper-gaps/cpgsv17_mpdo_theorem_4_9_implication_label.tex.
% **Scope restriction (single supplied factorization):** the theorem starts
% from one neighboring trace factorization of its full input tensor, including
% the rank-one identity tr(eta_{k,h}) = a_k b_h. Theorem 4.9(iv) supplies this
% data only for each BNT representative and leaves raw repeated-copy weights
% unconstrained. The printed (iv)=>(v) is therefore not an all-sector extension
% of this theorem and is refuted. Documented in
% docs/paper-gaps/cpgsv17_pf_rank_one.tex.
\begin{theorem}[Conditional two-site blocking theorem from neighboring trace factorization]
    \label{thm:mpdo_blocked_original_tensor_is_rfp_via_ts}
    \lean{MPOTensor.PhysicalSectorFactorization.NeighboringTraceFactorization.blockTwo_isRFPViaTS}
    \leanok
    \uses{def:is_rfp_via_ts, def:mpdo_two_site_blocking,
        def:mpdo_neighboring_trace_factorization}
    Let $F$ be a physical-sector factorization of $\mathcal K$ whose
    neighboring operators satisfy
    Definition~\ref{def:mpdo_neighboring_trace_factorization}. Thus they are
    positive semidefinite and there are real families $(a_k)_k,(b_k)_k$ with
    $\tr(\eta_{k,h})=a_kb_h$ and $\sum_k a_kb_k=1$. Let $\mathcal K^{[2]}$ be
    the tensor obtained by blocking two adjacent sites of $\mathcal K$. Then
    $\mathcal K^{[2]}$ is a renormalization fixed point in the sense of
    Definition~\ref{def:is_rfp_via_ts}. This is the conditional channel
    construction of \cite[Proposition~C.7]{Cirac2016MPDO_arXiv} for one
    supplied factorization. It does not extend to the printed all-BNT
    implication \cite[Theorem~4.9, (iv)$\Rightarrow$(v)]
      {Cirac2016MPDO_arXiv}: condition~(iv) constrains each BNT representative
    but not its raw repeated-copy weights, and that implication is refuted by
    Theorem~\ref{thm:mpdo_theorem49_iv_to_v_counterexample}. After condition
    (ii) has forced equal copy weights and their common values have been
    absorbed, these one-factorization channels become inputs to the
    projector-controlled assembly used by the viable implication
    (ii)$\Rightarrow$(v).
\end{theorem}

\begin{proof}\leanok
    \uses{thm:mpdo_blocked_sector_coordinate_tensor_is_rfp_via_ts}
    Transport the refinement and coarse-graining channels for
    $\widehat{\mathcal K}^{[2]}$ through the two-site and four-site tensor
    powers of the physical isometry, as in
    (\ref{eq:mpdo_physical_isometry_transport}). These unitary congruences
    preserve complete positivity and trace, while the corresponding closure
    identities carry the two fixed-point equations back to
    $\mathcal K^{[2]}$.
\end{proof}

% **Missing statement; proof-path drift:** Positivity, SAL, and literal ZCL
% descend to every common-weight absorbed BNT representative.  The printed
% per-representative factorization remains #6775; a direct sector-coarsening or
% sector-mixing construction of the same blocked channel pairs is #6793; and
% their composition through the orthogonal outer physical sectors is #6632.
% The raw implication (iv)=>(v) is separately refuted.  Documented in
% docs/paper-gaps/cpgsv17_pf_rank_one.tex.
\begin{theorem}[Arbitrary virtual boundaries of canonical block sums]
    \label{thm:mpdo_canonical_block_sum_arbitrary_boundary}
    \lean{MPSTensor.trace_evalWord_toTensorFromBlocks_mul,
      MPSTensor.trace_evalWord_gauge_mul,
      MPSTensor.trace_evalWord_gauge_toTensorFromBlocks_mul}
    \leanok
    Let a tensor be gauge equivalent, by an invertible virtual matrix $G$, to
    a finite block sum with component tensors $A_s$ and weights $\mu_s$. For
    an arbitrary virtual boundary $X$ and an arbitrary word $w$, its closure
    is the sum of the component closures with boundaries given by the diagonal
    blocks of $G^{-1}XG$:
    \begin{align}
      \tr\bigl(B^wX\bigr)
      &=\sum_s \mu_s^{|w|}\tr\bigl(A_s^w(G^{-1}XG)_{ss}\bigr).
      \notag
    \end{align}
    The conjugated boundary $G^{-1}XG$ is independent of the word length.
    This is the arbitrary-boundary form of the block calculation used in
    \cite[lines~186--195 and Appendix~C.2, lines~1646--1665 and 1810--1825]
      {Cirac2016MPDO_arXiv}.
\end{theorem}

\begin{proof}\leanok
    A word of length $|w|$ in the block-sum tensor is block diagonal. Its
    $s$-th diagonal block is $\mu_s^{|w|}A_s^w$. On multiplying by an
    arbitrary boundary matrix, only the diagonal boundary blocks contribute
    to the trace, which gives the asserted sum before the gauge change.

    Under the gauge $G$, the word $A^w$ is conjugated to $GA^wG^{-1}$.
    Cyclicity of the trace moves the two gauge factors to the boundary and
    replaces $X$ by $G^{-1}XG$. Combining these two identities proves the
    formula. Since this conjugation depends only on $G$ and $X$, the
    transformed boundary is the same at every word length.
\end{proof}

\begin{theorem}[Absorbed-representative boundary decomposition]
    \label{thm:mpdo_absorbed_representative_boundary_decomposition}
    \lean{MPOTensor.physCloseN_eq_sum_flatBasis_of_gauge_toTensor,
      MPOTensor.physCloseN_eq_sum_commonWeightAbsorbedBasis_of_gauge_toTensor}
    \leanok
    \uses{thm:mpdo_canonical_block_sum_arbitrary_boundary,
      def:mpdo_common_sector_weight_absorption, def:phys_close}
    % The basis blocks are an arbitrary family: the linked declarations take
    % a generic sector decomposition carrying no canonical-form or
    % basis-of-normal-tensors witness.
    Let the complete tensor $B$ be exactly gauge equivalent, by an invertible
    virtual matrix $G$, to a sector decomposition: an arbitrary finite family
    of basis blocks $A_j$, $j=1,\ldots,g$, each carrying a copy multiplicity
    $r_j$ and nonzero weights $\mu_{j,q}$. Suppose that the
    repeated-copy weights over each basis block are equal, with no equality
    required between different basis blocks, and absorb this common weight
    into its representative $\mathcal K_j$ as in
    Definition~\ref{def:mpdo_common_sector_weight_absorption}. Then, for every
    length $N$ and every virtual boundary $X$, the physical closure of the
    complete tensor is the sum of the absorbed-representative closures, with
    the boundary of a representative given by the sum of the diagonal copy
    blocks of $G^{-1}XG$ lying over it:
    \begin{align}
        B_N(X)
         & =
        \sum_{j=1}^g (\mathcal K_j)_N\Bigl(\sum_{q=1}^{r_j}
          (G^{-1}XG)_{(j,q),(j,q)}\Bigr).
        \notag
    \end{align}
    Neither the conjugation $X\mapsto G^{-1}XG$ nor the subsequent sum of
    diagonal copy blocks depends on $N$. In particular, the same
    representative boundary map applies at the two lengths entering the two
    channel equations of Definition~\ref{def:is_rfp_via_ts}. This is the
    arbitrary-boundary form of the block calculation used in
    \cite[lines~186--195 and Appendix~C.2, lines~1646--1665 and 1810--1825]
      {Cirac2016MPDO_arXiv}.
\end{theorem}

\begin{proof}\leanok
    Apply Theorem~\ref{thm:mpdo_canonical_block_sum_arbitrary_boundary} to the
    block sum indexed by the individual copies. Reindex this finite sum by a
    representative $j$ and a copy label $q$. For fixed $j$, all copy weights
    are equal, so their common $N$-th power may be absorbed into the length
    $N$ word of the representative tensor. The remaining dependence on $q$
    is linear in the diagonal boundary block. Summing over $q$ therefore
    replaces these blocks by their sum, which is the stated representative
    boundary. Neither the conjugation $X\mapsto G^{-1}XG$ nor the subsequent
    sum of diagonal copy blocks depends on $N$.
\end{proof}

\begin{theorem}[Conditional composition of absorbed BNT channels]
    \label{thm:mpdo_absorbed_bnt_channel_assembly}
    \lean{MPOTensor.blockedOneChainEquiv,
      MPOTensor.blockedTwoChainEquiv,
      MPOTensor.equivReindexMap_symm_apply_self,
      MPOTensor.equivReindexMap_physCloseN_two_eq_physClose1_blockTwo,
      MPOTensor.equivReindexMap_physCloseN_four_eq_physClose2_blockTwo,
      MPOTensor.blockedFirstSiteProjection,
      MPOTensor.blockedFirstSiteProjection_eq_reindex_kronecker_one,
      MPOTensor.firstSiteMatrix_mul_physCloseN_of_mul_physicalSlice,
      MPOTensor.firstSiteMatrix_mul_physCloseN_eq_zero_of_mul_physicalSlice_eq_zero,
      MPOTensor.physCloseN_mul_firstSiteMatrix_of_physicalSlice_mul,
      MPOTensor.singleKrausMap_firstSiteMatrix_physCloseN_of_twoSidedSupport,
      MPOTensor.firstSiteMatrix_isStarProjection,
      MPOTensor.sum_firstSiteMatrix_eq_one,
      MPOTensor.singleKrausMap_firstSiteMatrix_physCloseN_eq_ite,
      MPOTensor.coarseningMapInChainCoordinates,
      MPOTensor.refinementMapInChainCoordinates,
      MPOTensor.coarseningMapInChainCoordinates_isKrausCPTP,
      MPOTensor.refinementMapInChainCoordinates_isKrausCPTP,
      MPOTensor.coarseningMapInChainCoordinates_physCloseN,
      MPOTensor.refinementMapInChainCoordinates_physCloseN,
      MPOTensor.exists_chainCoordinateRFP_of_projectiveSectorDecomposition,
      MPOTensor.blockTwo_isRFPViaTS_of_chainCoordinateRFP,
      MPOTensor.blockTwo_isRFPViaTS_of_commonWeightAbsorbedBNTChannels}
    \leanok
    \uses{def:is_rfp_via_ts, def:mpdo_two_site_blocking,
        def:mpdo_common_sector_weight_absorption,
        thm:mpdo_absorbed_representative_boundary_decomposition,
        thm:mpdo_projective_resolution_channel_cptp}
    Let $K$ be exactly gauge equivalent, by an invertible virtual matrix $G$,
    to a canonical BNT sum. Suppose that the repeated-copy weights agree
    within each fixed representative and have been absorbed into tensors
    $(\mathcal K_s)_s$. Let $(P_s)_s$ be orthogonal projections on the
    one-site physical space such that
    \begin{align}
      \sum_s P_s&=I, & P_sP_t&=0\quad(s\ne t).
      \notag
    \end{align}
    Suppose also that, for every $s$ and every pair of virtual indices
    $\beta,\alpha$, the corresponding physical slice satisfies
    \begin{align}
      P_s\mathcal K_s[\beta,\alpha]
        &=\mathcal K_s[\beta,\alpha]
        =\mathcal K_s[\beta,\alpha]P_s .
      \notag
    \end{align}
    If every blocked representative $\mathcal K_s^{[2]}$ is a
    renormalization fixed point via trace-preserving completely positive maps,
    then $K^{[2]}$ is also such a renormalization fixed point.

    The global maps first measure the outer sector on the first original site
    of the blocked input. On one blocked site the corresponding projection is
    $P_s\otimes I$; on two blocked sites it acts on the first original site
    and as the identity on the remaining three. For every virtual boundary
    $X$, the boundary supplied to sector $s$ is the sum of the diagonal copy
    blocks of $G^{-1}XG$ over that representative. The same boundary is used
    in the length-two and length-four identities. This is the conditional
    composition underlying \cite[Appendix~C.2, lines~1646--1665 and
    1810--1825]{Cirac2016MPDO_arXiv}; it is not the printed implication
    (iv)$\Rightarrow$(v) of Theorem~4.9.
\end{theorem}

\begin{proof}\leanok
    \uses{thm:mpdo_absorbed_representative_boundary_decomposition,
        thm:mpdo_projective_resolution_channel_cptp}
    Write $\mathcal C_N(A,X)$ for the length-$N$ closure of $A$ with boundary
    $X$, and write $B_s(X)$ for the boundary supplied to representative $s$.
    At both relevant lengths the boundary decomposition gives
    \begin{align}
      \mathcal C_N(K,X)
        &=\sum_t \mathcal C_N(\mathcal K_t,B_t(X)),
        &N&\in\{2,4\}.
      \notag
    \end{align}
    Let $\Phi_{s,N}$ be compression by $P_s$ on the first site and by the
    identity on the remaining $N$ sites. Two-sided physical support and
    $P_sP_t=0$ for $s\ne t$ give the selection identity
    \begin{align}
      \Phi_{s,N}\bigl(\mathcal C_{N+1}(\mathcal K_t,B_t(X))\bigr)
        &=
        \begin{cases}
          \mathcal C_{N+1}(\mathcal K_t,B_t(X)),&s=t,\\
          0,&s\ne t.
        \end{cases}
      \notag
    \end{align}
    Apply the representative channel to the surviving summand and sum over
    $s$. The projective resolution makes both resulting maps trace-preserving
    and completely positive. The boundary decomposition then identifies their
    outputs with the length-two and length-four closures of $K$ carrying the
    original boundary.
\end{proof}

\begin{theorem}[Completion outside the absorbed BNT supports]
    \label{thm:mpdo_absorbed_bnt_channel_support_completion}
    \lean{MPOTensor.blockTwo_isRFPViaTS_of_commonWeightAbsorbedBNTChannels_of_orthogonalSupports}
    \leanok
    \uses{thm:mpdo_absorbed_representative_boundary_decomposition,
        thm:mpdo_absorbed_bnt_channel_assembly,
        def:is_rfp_via_ts}
    Retain the hypotheses of
    Theorem~\ref{thm:mpdo_absorbed_bnt_channel_assembly}, except that the
    pairwise orthogonal outer support projections are not required to sum to
    the identity. If every blocked absorbed representative is a
    renormalization fixed point via trace-preserving completely positive maps,
    then the blocked complete tensor is also such a fixed point.
\end{theorem}

\begin{proof}\leanok
    \uses{thm:mpdo_absorbed_representative_boundary_decomposition,
        thm:mpdo_absorbed_bnt_channel_assembly,
        thm:matrix_support_completion_cptp,
        lem:mpdo_compressed_cp_sum,
        lem:mpdo_projector_controlled_trace_selection}
    If the one-site physical space has dimension zero, all of its endomorphisms
    are equal. Hence $\sum_sP_s=I$, and
    Theorem~\ref{thm:mpdo_absorbed_bnt_channel_assembly} applies directly.
    Assume from now on that the physical dimension is positive. Fixed basis
    vectors witness that the two output spaces are nonempty; normalizing
    their identity matrices gives maximally mixed density matrices
    $\rho_2\in\MN{q^2}$ and $\rho_4\in\MN{q^4}$.

    Let $Q=\sum_sP_s$, an orthogonal projection by pairwise orthogonality.
    For $N\in\{2,4\}$ write $Q_N=Q\otimes\Id^{\otimes(N-1)}$ for the lift of
    $Q$ to the first original site of an $N$-site chain, and likewise
    $P_{s,N}=P_s\otimes\Id^{\otimes(N-1)}$, so that $Q_N=\sum_sP_{s,N}$.
    Two-sided support of each absorbed representative carries over to the
    length-two and length-four closures of the complete tensor $K$,
    \begin{align}
        Q_2K_2(X)Q_2&=K_2(X), & Q_4K_4(X)Q_4&=K_4(X),
        \label{eq:mpdo_support_completion_supported}
    \end{align}
    for every virtual boundary $X$.

    By hypothesis every $\mathcal K_s^{[2]}$ carries trace-preserving
    completely positive maps $\mathcal S_s,\mathcal T_s$ witnessing its
    fixed-point property. Sum these channels after compression by the
    matching first-site projection,
    $\mathcal S_{\mathrm{act}}(Y)=\sum_s\mathcal S_s(P_{s,4}YP_{s,4})$ and
    $\mathcal T_{\mathrm{act}}(Z)=\sum_s\mathcal T_s(P_{s,2}ZP_{s,2})$. Each
    sum is completely positive, and pairwise orthogonality of the $P_s$
    makes it preserve exactly the trace selected by $Q$,
    \begin{align}
        \tr\bigl(\mathcal S_{\mathrm{act}}(Y)\bigr)&=\tr(Q_4Y), &
        \tr\bigl(\mathcal T_{\mathrm{act}}(Z)\bigr)&=\tr(Q_2Z).
        \notag
    \end{align}

    Since $Q_N$ is Hermitian and idempotent, applying
    Theorem~\ref{thm:matrix_support_completion_cptp} to each active sum
    with $P=Q_N$ completes it on the orthogonal input complement,
    \begin{align}
        \mathcal S_{\mathrm{glob}}(Y)&=\mathcal S_{\mathrm{act}}(Y)
          +\tr\bigl((I-Q_4)Y\bigr)\rho_2, &
        \mathcal T_{\mathrm{glob}}(Z)&=\mathcal T_{\mathrm{act}}(Z)
          +\tr\bigl((I-Q_2)Z\bigr)\rho_4,
        \notag
    \end{align}
    and both global maps are trace-preserving and completely positive. By
    (\ref{eq:matrix_support_completion_agreement}) and
    (\ref{eq:mpdo_support_completion_supported}),
    $\mathcal S_{\mathrm{glob}}(K_4(X))=\mathcal S_{\mathrm{act}}(K_4(X))$ and
    $\mathcal T_{\mathrm{glob}}(K_2(X))=\mathcal T_{\mathrm{act}}(K_2(X))$ for
    every virtual boundary $X$, so the two channel equations for
    $\mathcal S_{\mathrm{glob}},\mathcal T_{\mathrm{glob}}$ reduce to those of
    $\mathcal S_{\mathrm{act}},\mathcal T_{\mathrm{act}}$ and hence, by
    construction of the active sums, to the representative equations.
\end{proof}

\begin{theorem}[Blocking channels after common-weight absorption]
    \label{thm:mpdo_sal_zcl_blocking_channels}
    \notready
    \uses{def:mpdo_simple_cf_tensor, def:sector_bnt_cf,
        def:is_sal, def:mpo_physical_trace_transfer, def:is_rfp_via_ts,
        def:mpdo_two_site_blocking}
    Let $K$ be a simple tensor in block-injective canonical form (biCF),
    equipped with a basis-of-normal-tensors decomposition and satisfying the
    strong area law and literal zero correlation length,
    \begin{align}
        \mathcal T_K^2&=\mathcal T_K.
        \notag
    \end{align}
    Let $M=K^{[2]}$ be its two-site blocking. Then
    there are trace-preserving completely positive maps in both directions
    between the one-site and two-site closures of $M$. Equivalently, $M$ is a
    renormalization fixed point via these two maps. This is implication
    (ii)$\Rightarrow$(v) of
    \cite[Theorem~4.9 and Appendix~C.2, lines~1810--1825]
      {Cirac2016MPDO_arXiv}. The separate standing biCF and BNT hypotheses are
    imposed at lines~849--850 and restated at line~1628 of the same source.
    The source first uses zero correlation length at lines~1646--1665 to make
    repeated-copy weights independent of the copy label and absorbs their
    common value into each representative. The proof at lines~1815--1816 then
    projects onto the mutually orthogonal outer physical sectors and invokes
    the one-representative channel construction. Thus a faithful proof must
    supply one blocked channel pair for every absorbed representative. Their
    combination through the outer physical projections, including the
    trace-preserving completion outside their sum, is
    Theorem~\ref{thm:mpdo_absorbed_bnt_channel_support_completion}.
\end{theorem}

% Correct equation, not a figure: the cited source draws the two-to-three
% channel MPDO_TSKKK.png after explicitly suppressing U, but it does not draw
% this tensor-power transport square.  The transport theorem is recorded in
% docs/paper-gaps/cpgsv17_mpdo_blocked_rfp_physical_transport.tex.
\begin{align}
    U^{\otimes4}\widehat{\mathcal T}^{\,2}
      (U^\dagger)^{\otimes2}
    &=\mathcal T,
    \label{eq:mpdo_physical_isometry_transport}\\
    U^{\otimes2}\widehat{\mathcal S}^{\,2}
      (U^\dagger)^{\otimes4}
    &=\mathcal S.
    \notag
\end{align}
Conjugation by $U^{\otimes n}$ and $(U^\dagger)^{\otimes n}$ carries the
sector-coordinate channels to the ordinary-coordinate channels. This combines
Proposition~C.7 with
\cite[Appendix~C.2, lines~1510--1563 and 1821--1825]{Cirac2016MPDO_arXiv}.

Tracing the ket and bra indices of a sector tensor gives the bond vector
$|l_k)$ or the bond functional $(r_k|$. Hence a neighboring pair gives
$T_{k,h}=(r_k|l_h)=\tr(\eta_{k,h})$, as in
\cite[Appendix~C.2, lines~1474--1481]{Cirac2016MPDO_arXiv}.
